% 『物理を理解するための数学 問題集』第8巻（第42回編集委員会）
% 本文・一覧・こたえは mondai/build/mk_book.js が書き出す（mq8-body・mq8-list・mq8-ans）。
\documentclass[12pt]{ltjsarticle}
\def\bkfontpath{fonts/}
\usepackage[vol=8]{bk-b5}
\usepackage{bk-matter}
\usepackage{bk-honbun}
\usepackage{bk-volume}
\usepackage{bk-pict1}
\usepackage{bk-mondai}
\hypersetup{pdftitle={物理を理解するための数学 問題集　第8巻　空間に広がる量}, pdfauthor={「物理を理解するための数学」編集委員会}}
\begin{document}
% 表紙のわざの見本・使い方の「たとえば」（第42回。事務局の仮：わざができたら書く組・まとめ役が決め直す）
\def\mqcoverwaza{0/8-10,1/8-27,2/8-69}
\def\mqhowtoex{たとえば $x^3y^2$ を $x$ で偏微分するときは、$y^2$ を数とみて止めておけば、$3x^2y^2$ とすぐわかります。}
\def\mqhowtowaza{8-10}
\def\mqhowtonote{ステップ3とわなの図は、形を示すものです。図の矢印の長さや目もりを読んで答えを出すのではなく、わざを使って計算で求めましょう。}
%% 問題集の前付・こたえ・奥付（全巻共通。第29回）
\makeatletter
% 表紙
\newcommand\mqcover[2]{%
  \thispagestyle{empty}\bkdecoration{false}\bkfolio{false}%
  \begin{tikzpicture}[remember picture, overlay, x=1mm, y=1mm]
    \fill[bkTheme] ([xshift=-3mm,yshift=3mm]current page.north west) rectangle ([xshift=3mm,yshift=-3mm]current page.south east);
    \fill[white] ([xshift=14mm,yshift=-58mm]current page.north west) rectangle ([xshift=-14mm,yshift=-150mm]current page.north east);
    \node[anchor=north west, font=\sffamily\bfseries\large, text=white] at ([xshift=16mm,yshift=-18mm]current page.north west) {\R{物理}{ぶつ|り}を\R{理解}{り|かい}するための\R{数学}{すう|がく}};
    \node[anchor=north west, font=\sffamily\bfseries\fontsize{30}{36}\selectfont, text=white] at ([xshift=16mm,yshift=-30mm]current page.north west) {\R{問題集}{もん|だい|しゅう}};
    \node[anchor=north west, font=\sffamily\bfseries\Large, text=bkDeep] at ([xshift=20mm,yshift=-64mm]current page.north west) {\R{第}{だい}#1\R{巻}{かん}};
    \node[anchor=north west, font=\bfseries\fontsize{30}{38}\selectfont, text=bkInk] at ([xshift=20mm,yshift=-80mm]current page.north west) {#2};
    \node[anchor=north west, font=\sffamily\bfseries\large, text=bkDeep, text width=140mm] at ([xshift=20mm,yshift=-110mm]current page.north west) {けいさんを らくに する 「わざ」を\\ \R{見}{み}わける・\R{作}{つく}る};
    \foreach \x/\t in {0/1-02,1/1-37,2/1-43} \node[fill=white, text=bkDeep, rounded corners=3mm, font=\sffamily\bfseries\Large, inner xsep=4mm, inner ysep=2.5mm]
      at ([xshift={30mm+\x*42mm},yshift=-178mm]current page.north west) {わざ\hspace{0.3em}\t};
    \node[anchor=south west, font=\sffamily\small, text=white] at ([xshift=16mm,yshift=14mm]current page.south west) {「\R{物理}{ぶつ|り}を\R{理解}{り|かい}するための\R{数学}{すう|がく}」\R{編集}{へん|しゅう}\R{委員会}{い|いん|かい}};
  \end{tikzpicture}\null\clearpage}
% 使い方（見開き：2〜3ページ）
\newcommand\mqhowto{%
  \bkdecoration{true}\bkfolio{true}\bksetrunning{}{\R{問題集}{もん|だい|しゅう}\hspace{0.5em}\R{第}{だい}\bkvol\R{巻}{かん}}{この ほんの つかいかた}%
  \noindent{\sffamily\bfseries\Large\color{bkDeep}この ほんの つかいかた}\par\vspace{2mm}
  \noindent この ほんは、 けいさんを らくに する 「わざ」を あつめた もんだいしゅうです。 たとえば 9 × 7 は、 10 × 7 ＝ 70 から 7を ひけば、 すぐ 63と わかります。 こういう わざを、 1つずつ ばんごうを つけて ならべました。\par\vspace{2mm}
  \begin{tcolorbox}[enhanced, colback=bkTheme!5, colframe=bkTheme, boxrule=1pt, arc=2mm, left=2mm, right=2mm, top=1mm, bottom=1mm,
    fonttitle=\sffamily\bfseries\small, coltitle=white, colbacktitle=bkTheme, title={1つの わざは、 \R{見}{み}ひらき 2ページ}]
  \textbf{\R{左}{ひだり}の ページ}\par
  \mqwazachipsmall{わざ 1-37}\ わざの ばんごうと なまえ\par
  \R{見}{み}わける めじるし：どんな もんだいの ときに つかえるか\par
  ふつうに とくと ／ らくな かたちに すると：くらべて みよう\par
  どうして らくに なるの？・つかえない とき\par
  ステップ1：そのまま つかう\par\vspace{1mm}
  \textbf{\R{右}{みぎ}の ページ}\par
  ステップ2：すこし かたちを かえる\par
  ステップ3：\R{見}{み}かけが ちがう\par
  わなに ちゅうい：にて いるけど、 らくに ならない もんだい
  \end{tcolorbox}
  \par\vspace{2mm}\noindent ぶの さいごには、 いろいろな わざが まざった 「\R{見}{み}わけよう」と、 じぶんで もんだいを つくる 「\R{作}{つく}って みよう」が あります。\par
  \clearpage
  \noindent{\sffamily\bfseries\Large\color{bkDeep}こたえと プリント}\par\vspace{2mm}
  \noindent こたえは、 ほんの うしろに あります。 ［\ ］の なかは、 らくな かたちです。 こたえが あって いても、 らくな かたちを くらべて みよう。\par\vspace{3mm}
  \begin{tcolorbox}[enhanced, colback=white, colframe=bkDeep, boxrule=1pt, arc=2mm, left=2mm, right=2mm, top=1mm, bottom=1mm,
    fonttitle=\sffamily\bfseries\small, coltitle=white, colbacktitle=bkDeep, title={プリントを \R{作}{つく}る アプリ}]
  わざの ばんごう（たとえば 1-37）を、 プリントを \R{作}{つく}る アプリに \R{入}{い}れると、 その わざの もんだいが、 なんまいでも できます。 プリントの \R{上}{うえ}の 「たねの ばんごう」が おなじなら、 おなじ プリントが もう いちど できます。
  \end{tcolorbox}\par\vspace{3mm}
  \noindent{\sffamily\bfseries\color{bkDeep}おうちの \R{方}{かた}へ}\par
  {\small この\R{問題集}{もん|だい|しゅう}は、「\R{物理}{ぶつ|り}を\R{理解}{り|かい}するための\R{数学}{すう|がく}」\R{第}{だい}1\R{巻}{かん}に\R{対応}{たい|おう}しています。わざの\R{見開}{み|ひら}きには、\R{教科書}{きょう|か|しょ}のどの\R{章}{しょう}で\R{学}{まな}ぶかを\R{示}{しめ}しました（\mqtextbook{1-13} のような\R{札}{ふだ}）。★★のわざは、\R{第}{だい}1\R{巻}{かん}の\R{少}{すこ}し\R{先}{さき}（2けた×1けたなど）の\R{内容}{ない|よう}です。わざを\R{暗記}{あん|き}するより、「どうして\R{楽}{らく}になるのか」と「\R{使}{つか}えないとき」を\R{一緒}{いっ|しょ}に\R{確}{たし}かめてください。}\par
  \clearpage}
% わざの一覧
\newcommand\mqlistpages{%
  \bksetrunning{}{\R{問題集}{もん|だい|しゅう}\hspace{0.5em}\R{第}{だい}\bkvol\R{巻}{かん}}{わざの いちらん}%
  \noindent{\sffamily\bfseries\Large\color{bkDeep}わざの いちらん}\par\vspace{1mm}
  \noindent{\small ばんごうを アプリに \R{入}{い}れると、 その わざの プリントが できます。 \mqstars{2}は すこし むずかしい わざ。}\par\vspace{1mm}
  \input{mq\bkvol-list}\clearpage}
\newcommand\mqfrontmatter[2]{\mqcover{#1}{#2}\mqhowto\mqlistpages}
% こたえ
\newcommand\mqanswers[1]{%
  \clearpage\ifodd\value{page}\else\mqmemopage\fi\typeout{MQSPREAD \thepage}%
  \bksetrunning{}{\R{問題集}{もん|だい|しゅう}\hspace{0.5em}\R{第}{だい}\bkvol\R{巻}{かん}}{こたえ}%
  \phantomsection\pdfbookmark[0]{こたえ}{mqans}%
  \noindent{\sffamily\bfseries\Large\color{bkDeep}こたえ}\hspace{1em}{\small ［\ ］の なかは らくな かたち}\par\vspace{1mm}
  \begin{multicols}{2}\raggedright\footnotesize\setlength{\parskip}{0pt}
  #1
  \end{multicols}}
% 奥付
\newcommand\mqcolophon[3]{%
  \clearpage\bkdecoration{false}%
  \null\vfill
  \noindent\begin{tabular}{@{}ll@{}}
  \multicolumn{2}{@{}l}{\sffamily\bfseries\large \R{物理}{ぶつ|り}を\R{理解}{り|かい}するための\R{数学}{すう|がく}\ \R{問題集}{もん|だい|しゅう}\ \R{第}{だい}#1\R{巻}{かん}\ #2}\\[2mm]
  {\small \R{版}{はん}} & {\small #3}\\
  {\small \R{編}{へん}} & {\small 「\R{物理}{ぶつ|り}を\R{理解}{り|かい}するための\R{数学}{すう|がく}」\R{編集}{へん|しゅう}\R{委員会}{い|いん|かい}}\\
  {\small \R{問題}{もん|だい}} & {\small \R{種}{たね}の\R{決}{き}まった\R{乱数}{らん|すう}で\R{作}{つく}り、\R{検算}{けん|ざん}のプログラムで\R{確}{たし}かめた}\\
  \end{tabular}\par}
\makeatother

\mqfrontmatter{8}{空間に広がる量}
\input{mq8-body}
\mqanswers{\input{mq8-ans}}
\mqcolophon{8}{空間に広がる量}{第1版 Ver1　2026年11月}
\end{document}
